\clearpage\begingroup\NoSourceHead\SourceLeaf{515b}
\renewcommand{\thefootnote}{\arabic{footnote})}\setcounter{footnote}{7}
\newcommand{\BolzaEqualAboveLess}{\mathrel{\vcenter{\offinterlineskip\hbox{$\scriptstyle=$}\kern-.15ex\hbox{$<$}}}}
\ArithmeticHead{58}{L}{BOLZA, GAUSS UND DIE VARIATIONSRECHNUNG.}
\noindent\hspace*{\parindent}c) Verwandte Sätze und weitere Verallgemeinerungen.
\par\noindent\hspace*{\parindent}8. An den \textsc{Kneser}schen Transversalensatz, der als eine unmittelbare Verallgemeinerung
des \textsc{Gauss}schen Satzes erscheint, schliessen sich als »verwandte«
Sätze in dem oben definierten Sinn der Satz über \textls[100]{Parallelflächen} an, der
sich auf das Extremum des Integrals
\begin{equation*}\tag{18}
\int\sqrt{x'^2+y'^2+z'^2}\,dt
\end{equation*}
bezieht, sowie der Satz von \textsc{Thomson} und \textsc{Tait}\footnote{W. \textsc{Thomson} und P. G. \textsc{Tait}, \textit{Treatise on Natural Philosophy}, Oxford 1867. Bd. 1, S. 244;
deutsche Übersetzung von \textsc{Helmholtz} und \textsc{Wertheim} unter dem Titel: \textit{Handbuch der theoretischen Physik},
Bd. I, Teil I, Braunschweig 1871, S. 272.} über Oberflächen gleicher
\clearpage\SourceLeaf{516}
\ArithmeticHead{59}{R}{DER TRANSVERSALENSATZ.}
\noindent Wirkung und der Satz von \textsc{Malus}\footnote{Zuerst von W. R. \textsc{Hamilton} ausgesprochen, \textit{Theory of systems of rays}, Transactions of the Royal
Irish Academy, Bd. 15 (1828), Einleitung Nr. 106. Vgl. auch H. \textsc{Weber}, \textit{Ueber den Satz von Malus für
krummlinige Lichtstrahlen}, Rendiconti del Circolo Matematico di Palermo, Bd. 29 (1910), S. 404.} für krummlinige Lichtstrahlen in einem
Medium von stetig veränderlichem Brechungsindex, bei denen das vorige Integral
durch das etwas allgemeinere
\[
\int n(x,y,z)\sqrt{x'^2+y'^2+z'^2}\,dt
\]
ersetzt ist.

Diese Sätze sind mit dem \textsc{Gauss}schen Satz zusammen als Spezialfälle in
dem Transversalensatz für das Integral
\begin{equation*}\tag{19}
J=\int F(x_1,x_2,\ldots,x_n,x'_1,x'_2,\ldots,x'_n)\,dt\equiv\int F(x,x')\,dt
\end{equation*}
enthalten, wo wieder $F$ positiv-homogen von der ersten Dimension in
$x'_1,x'_2,\ldots,x'_n$ ist. Hier handelt es sich dann im allgemeinsten Fall um eine
$m$-parametrige Extremalenschar dieses Integrals $(1\BolzaEqualAboveLess m\BolzaEqualAboveLess n-1)$ mit Parametern
$v_1,v_2,\ldots,v_m$ und um eine $m$-dimensionale Mannigfaltigkeit
\[
\mathfrak C_0:x_k=\xi_k(v_1,v_2,\ldots,v_m),\qquad k=1,2,\ldots,n,
\]
welche die Extremalenschar transversal schneidet, d. h. für die im Schnittpunkt
mit der Extremale $v_1,v_2,\ldots,v_m$ die $m$ Gleichungen erfüllt sind
\[
\sum_{i=1}^n F_{x'_i}(x,x')\frac{\partial\xi_i}{\partial v_h}=0,\qquad h=1,2,\ldots,m.
\]
Dementsprechend treten jetzt an Stelle der einen Funktion $S(u,v)$ die $m$
Funktionen
\begin{equation*}\tag{20}
S_h(u,v_1,v_2,\ldots,v_m)=\sum_{i=1}^n F_{x'_i}(X,X')\frac{\partial X_i}{\partial v_h},\qquad h=1,2,\ldots,m,
\end{equation*}
wenn die gegebene Extremalenschar durch die Gleichungen
\[
x_k=X_k(u,v_1,v_2,\ldots,v_m),\qquad k=1,2,\ldots,n
\]
dargestellt ist und $u$ die analoge Bedeutung hat wie oben, sodass also:
\[
F(X,X')=1.
\]

Auch hier gilt dann die \textsc{Gauss}sche Fundamentalgleichung
\begin{equation*}\tag{21}
\frac{\partial S_h}{\partial u}=0,\qquad h=1,2,\ldots,m,
\end{equation*}
woraus der Transversalensatz genau so folgt wie im Fall $n=2$.
\par\smallskip\noindent\hfill$8^*$\par
\clearpage\SourceLeaf{517}
\ArithmeticHead{60}{L}{BOLZA, GAUSS UND DIE VARIATIONSRECHNUNG.}
\noindent\hspace*{\parindent}Ein weiterer »verwandter« Satz ist der Satz von \textsc{Malus} und \textsc{Dupin}\footnote{Über die Geschichte des Satzes vergleiche G. \textsc{Darboux}, \textit{Leçons sur la théorie générale des surfaces},
Bd. II (Paris 1889), S. 280. Die erste Arbeit von \textsc{Malus} stammt danach aus dem Jahr 1808.} über
reflektierte und gebrochene geradlinige Strahlensysteme mit dem Zusatz von
\textsc{Hamilton}\footnote{loc. cit. 9), Art. 13.}, dass für den Fall von Reflexionen die Länge $\Sigma\rho$ des polygonalen
Weges eines Lichtstrahls zwischen irgend zwei Orthogonalflächen konstant ist,
während für den Fall von Brechungen $\Sigma\rho$ durch $\Sigma n\rho$ zu ersetzen ist, wenn $n$
den Brechungsindex des betreffenden Mediums bedeutet.
\par\clearpage\endgroup
